\documentclass[a4paper,12pt]{article}

\usepackage[T2A]{fontenc}
\usepackage[utf8]{inputenc}
\usepackage[russian]{babel}
\usepackage{amsmath,amssymb,amsthm,mathtools}
\usepackage{helvet}
\renewcommand{\familydefault}{\sfdefault}
\usepackage{icomma}
\usepackage{geometry}
\geometry{left=20mm,right=20mm,top=20mm,bottom=22mm}
\usepackage{microtype}
\usepackage{tikz}
\usetikzlibrary{calc,arrows.meta,positioning,shapes.geometric}
\usepackage{tkz-euclide}
\usepackage{pgfplots}
\pgfplotsset{compat=1.18}
\usepackage{tabularx,booktabs,longtable}
\usepackage{enumitem}
\usepackage{hyperref}
\usepackage{parskip}
\usepackage{fancyhdr}
\usepackage{setspace}
\usepackage{xcolor}

\definecolor{accent}{HTML}{008080}
\definecolor{darkaccent}{HTML}{004D4D}
\definecolor{textgray}{HTML}{333333}
\definecolor{softgray}{HTML}{F3F5F5}

\hypersetup{colorlinks=true,linkcolor=darkaccent,urlcolor=darkaccent}
\onehalfspacing
\setlength{\parindent}{0pt}
\setlist[itemize]{leftmargin=6mm,itemsep=2pt,topsep=3pt}
\setlist[enumerate]{leftmargin=7mm,itemsep=4pt,topsep=4pt}
\renewcommand{\arraystretch}{1.3}
\setlength{\headheight}{25pt}

\pagestyle{fancy}
\fancyhf{}
\fancyhead[L]{\small\color{textgray}Графики функций: прямая и парабола}
\fancyhead[R]{\small\color{textgray}ОГЭ}
\fancyfoot[C]{\small\color{textgray}\thepage}
\renewcommand{\headrulewidth}{0.4pt}
\renewcommand{\headrule}{%
  \hbox to\headwidth{%
    \color{accent}\leaders\hrule height \headrulewidth\hfill
  }%
}

\newcommand{\sectionline}{%
  \par\noindent
  \color{accent}\rule{\linewidth}{0.5pt}
  \color{black}\par
}

\newcommand{\key}[1]{\textcolor{darkaccent}{\textbf{#1}}}

\tikzset{
  flowstart/.style={
    ellipse,
    draw=darkaccent,
    line width=0.9pt,
    minimum width=38mm,
    minimum height=10mm,
    align=center,
    inner sep=3pt
  },
  flowproc/.style={
    rectangle,
    rounded corners=2mm,
    draw=accent,
    line width=0.9pt,
    text width=104mm,
    minimum height=12mm,
    align=center,
    inner sep=5pt
  },
  flowcond/.style={
    diamond,
    aspect=2.6,
    draw=darkaccent,
    line width=0.9pt,
    text width=58mm,
    align=center,
    inner sep=2pt
  },
  flowend/.style={
    ellipse,
    draw=darkaccent,
    line width=0.9pt,
    minimum width=42mm,
    minimum height=10mm,
    align=center,
    inner sep=3pt
  },
  flowarrow/.style={
    -{Stealth[length=3mm,width=2mm]},
    line width=0.9pt,
    color=darkaccent
  },
  flowlabel/.style={
    font=\small,
    fill=white,
    inner sep=1.5pt,
    text=black
  }
}

\begin{document}

% =========================================================
% ТИТУЛЬНЫЙ ЛИСТ
% =========================================================

\begin{titlepage}
\thispagestyle{empty}

\begin{center}

\vspace*{24mm}

{\Large\textcolor{darkaccent}{\textbf{Чек-лист}}}

\vspace{8mm}

{\LARGE\textbf{Графики линейной и квадратной функций}}

\vspace{3mm}

{\Large\textbf{коэффициенты, пересечения, параметры}}

\vspace{15mm}

{\large Подготовка к ОГЭ по математике}

\vspace{18mm}

{\large \today}

\vfill

{\large Лёвин Артём Александрович}

\vspace{2mm}

{\large\textbf{эксклюзивно для Марины}}

\vspace{18mm}

\begin{tikzpicture}[x=1cm,y=1cm]
\draw[accent,line width=1.1pt]
  (0,0)
  .. controls (2.2,0.9) and (4.2,-0.9) .. (6.2,0)
  .. controls (8.1,0.8) and (10.1,-0.55) .. (12.3,0.15);
\end{tikzpicture}

\end{center}
\end{titlepage}

% =========================================================
% 1. ЛИНЕЙНАЯ ФУНКЦИЯ
% =========================================================

\section*{1. Линейная функция: что читать по графику}
\sectionline

Линейная функция имеет вид
\[
y=kx+b.
\]

\begin{tabularx}{\linewidth}{@{}>{\bfseries}p{0.18\linewidth}X@{}}
\toprule
Коэффициент & Что показывает \\
\midrule
$k$
&
Наклон прямой:
при $k>0$ функция возрастает,
при $k<0$ убывает,
при $k=0$ график горизонтален.
\\
$b$
&
Ординату точки пересечения с осью $Oy$:
при $x=0$ получаем $y=b$.
\\
\bottomrule
\end{tabularx}

\key{Абсцисса} --- координата $x$.

\key{Ордината} --- координата $y$.

\subsection*{Быстрая идентификация прямой}

Если нужно сопоставить график и формулу, удобно действовать в два этапа:

\begin{enumerate}
  \item по направлению прямой определить знак $k$;
  \item взять узловую точку графика с целыми координатами и проверить её подстановкой в оставшиеся формулы.
\end{enumerate}

\textbf{Пример.}
Точка $(1;-3)$ лежит на графике.
Для формулы $y=-3x$:
\[
-3=-3\cdot1,
\]
поэтому точка удовлетворяет формуле.

\subsection*{Особые прямые}

\begin{minipage}[t]{0.49\linewidth}
\vspace{0pt}

\begin{itemize}
  \item $y=c$ --- горизонтальная прямая, параллельная оси $Ox$;

  \item $x=c$ --- вертикальная прямая, параллельная оси $Oy$; это \textbf{не} график функции $y=f(x)$;

  \item $y=m$ --- семейство горизонтальных прямых: при изменении $m$ прямая перемещается вверх и вниз;

  \item $y=mx$ --- семейство прямых через $(0;0)$: при изменении $m$ прямая поворачивается вокруг начала координат.
\end{itemize}

\end{minipage}
\hfill
\begin{minipage}[t]{0.47\linewidth}
\vspace{0pt}
\centering

\begin{tikzpicture}
\begin{axis}[
  width=\linewidth,
  height=5.7cm,
  axis lines=middle,
  xmin=-4.5,
  xmax=4.5,
  ymin=-4.5,
  ymax=5.5,
  xlabel={$x$},
  ylabel={$y$},
  xtick={-4,-2,0,2,4},
  ytick={-4,-2,0,2,4},
  grid=both,
  minor tick num=1,
  unit vector ratio=1 1,
  samples=220,
  clip=false,
  every axis plot/.append style={line width=1pt}
]

\addplot[
  accent,
  domain=-4.5:4.5
]{2};

\addplot[
  darkaccent,
  domain=-2.2:2.2
]{1.5*x};

\node[
  anchor=west,
  text=accent,
  font=\small
] at (axis cs:2.25,2.35) {$y=2$};

\node[
  anchor=west,
  text=darkaccent,
  font=\small
] at (axis cs:1.15,3.25) {$y=1{,}5x$};

\end{axis}
\end{tikzpicture}

\end{minipage}

% =========================================================
% 2. ПЕРЕСЕЧЕНИЕ ГРАФИКОВ
% =========================================================

\clearpage

\section*{2. Точка пересечения двух графиков}
\sectionline

В общей точке двух графиков совпадают и абсцисса, и ордината.
Поэтому для
\[
y=f(x),
\qquad
y=g(x)
\]
условие пересечения имеет вид
\[
\boxed{f(x)=g(x)}.
\]

\textbf{Алгоритм вычисления:}

\begin{enumerate}
  \item приравнять правые части функций;
  \item решить полученное уравнение относительно $x$;
  \item подставить найденное $x$ в любую исходную функцию и найти $y$;
  \item при желании проверить $y$ по второй функции.
\end{enumerate}

\textbf{Пример.}
Найдём точку пересечения
\[
y_1=2x+5,
\qquad
y_2=-x+7.
\]

Приравниваем:
\[
2x+5=-x+7
\quad\Longrightarrow\quad
3x=2
\quad\Longrightarrow\quad
x=\frac23.
\]

Тогда
\[
y
=
2\cdot\frac23+5
=
\frac{19}{3}.
\]

Следовательно,
\[
\boxed{\left(\frac23;\frac{19}{3}\right)}.
\]

Проверка по второй функции:
\[
-\frac23+7=\frac{19}{3}.
\]

\textbf{Важно.}
Тот же принцип работает для любых двух графиков, например для прямой и параболы:
сначала решается уравнение
\[
f(x)=g(x).
\]

% =========================================================
% БЛОК-СХЕМА 1
% =========================================================

\clearpage
\thispagestyle{fancy}

\begin{center}
{\Large\textcolor{darkaccent}{\textbf{Алгоритм: точка пересечения двух графиков}}}
\end{center}

\vspace{10mm}

\begin{center}
\begin{tikzpicture}[node distance=10mm]

\node[flowstart] (start)
{Начало};

\node[flowproc,below=of start] (eq)
{Запишите условие пересечения: $f(x)=g(x)$};

\node[flowproc,below=of eq] (solve)
{Решите уравнение и найдите допустимые значения $x$};

\node[flowcond,below=13mm of solve] (roots)
{Получены значения $x$?};

\node[flowproc,below=14mm of roots] (ys)
{Для каждого $x$ вычислите $y=f(x)$};

\node[flowproc,below=of ys] (check)
{Проверьте при необходимости: $y=g(x)$};

\node[flowend,below=of check] (end)
{Запишите точку или точки $(x;y)$};

\node[flowend,right=26mm of roots] (none)
{Общих точек нет};

\draw[flowarrow] (start) -- (eq);
\draw[flowarrow] (eq) -- (solve);
\draw[flowarrow] (solve) -- (roots);

\draw[flowarrow]
  (roots)
  --
  node[flowlabel,right]{да}
  (ys);

\draw[flowarrow] (ys) -- (check);
\draw[flowarrow] (check) -- (end);

\draw[flowarrow]
  (roots.east)
  --
  node[flowlabel,above]{нет}
  (none.west);

\end{tikzpicture}
\end{center}

\clearpage

% =========================================================
% 3. КВАДРАТИЧНАЯ ФУНКЦИЯ
% =========================================================

\section*{3. Квадратичная функция и парабола}
\sectionline

Квадратичная функция имеет вид
\[
\boxed{y=ax^2+bx+c},
\qquad
a\ne0.
\]

Условие $a\ne0$ принципиально:
при $a=0$ квадратный член исчезает и получается линейная функция
\[
y=bx+c.
\]

\subsection*{Коэффициент $a$: направление ветвей}

\[
\begin{array}{ccl}
a>0
&
\Longrightarrow
&
\text{ветви направлены вверх},
\\[2mm]
a<0
&
\Longrightarrow
&
\text{ветви направлены вниз}.
\end{array}
\]

\subsection*{Вершина параболы и коэффициент $b$}

Абсцисса вершины вычисляется по формуле
\[
\boxed{x_V=-\frac{b}{2a}}.
\]

Отсюда полезная обратная запись:
\[
\boxed{b=-2ax_V}.
\]

Она позволяет определить знак $b$, если по графику видны знак $a$
и расположение вершины относительно оси $Oy$.

\begin{itemize}
  \item $x_V>0$: вершина находится справа от оси $Oy$;
  \item $x_V<0$: вершина находится слева от оси $Oy$;
  \item $x_V=0$: вершина лежит на оси $Oy$, тогда $b=0$.
\end{itemize}

\subsection*{Коэффициент $c$: пересечение с осью $Oy$}

При $x=0$:
\[
y=c.
\]

Значит, $c$ --- ордината точки пересечения параболы с осью $Oy$.

\begin{center}
\begin{tabularx}{0.88\linewidth}{@{}>{\bfseries}p{0.16\linewidth}X@{}}
\toprule
$a$ & ветви параболы \\
$b$ & положение вершины по формуле $x_V=-\dfrac{b}{2a}$ \\
$c$ & пересечение с осью $Oy$ \\
\bottomrule
\end{tabularx}
\end{center}

\begin{center}
\begin{tikzpicture}

\begin{axis}[
  width=0.86\linewidth,
  height=7.3cm,
  axis lines=middle,
  xmin=-4.5,
  xmax=5.5,
  ymin=-5.5,
  ymax=6.5,
  xlabel={$x$},
  ylabel={$y$},
  xtick={-4,-2,0,2,4},
  ytick={-4,-2,0,2,4,6},
  grid=both,
  minor tick num=1,
  unit vector ratio=1 1,
  samples=240,
  clip=false,
  every axis plot/.append style={line width=1.15pt}
]

\addplot[
  accent,
  domain=-3.1:5.1
]{(x-1)^2-4};

\addplot[
  darkaccent,
  domain=-3.2:4.2
]{-0.7*(x-0.5)^2+4};

\addplot[
  only marks,
  mark=*,
  mark size=1.8pt,
  accent
] coordinates {(1,-4)};

\addplot[
  only marks,
  mark=*,
  mark size=1.8pt,
  darkaccent
] coordinates {(0.5,4)};

\node[
  anchor=north west,
  text=accent
] at (axis cs:1.15,-4) {$V_1$};

\node[
  anchor=south west,
  text=darkaccent
] at (axis cs:0.65,4) {$V_2$};

\end{axis}
\end{tikzpicture}
\end{center}

% =========================================================
% 4. ЗНАКИ КОЭФФИЦИЕНТОВ
% =========================================================

\clearpage

\section*{4. Как определять знаки $a$, $b$, $c$ по графику}
\sectionline

Для параболы
\[
y=ax^2+bx+c
\]
удобно двигаться в фиксированном порядке.

\begin{enumerate}
  \item \key{Определите знак $a$} по направлению ветвей.

  \item \key{Определите знак $x_V$} по расположению вершины слева или справа от оси $Oy$.

  \item \key{Определите знак $b$} из соотношения
  \[
  b=-2ax_V.
  \]

  \item \key{Определите знак $c$} по точке пересечения с осью $Oy$.
\end{enumerate}

\textbf{Пример.}
Ветви направлены вверх, вершина находится справа от $Oy$,
а пересечение с $Oy$ расположено ниже оси $Ox$.

Тогда
\[
a>0,
\qquad
x_V>0,
\qquad
b=-2ax_V<0,
\qquad
c<0.
\]

\subsection*{Типичные ошибки}

\begin{itemize}
  \item путать знак $b$ со знаком абсциссы вершины: между ними есть ещё множитель $-2a$;

  \item забывать минус в формуле
  \[
  x_V=-\frac{b}{2a};
  \]

  \item определять $c$ по вершине: $c$ связан с точкой пересечения с осью $Oy$;

  \item считать $x=c$ функцией вида $y=f(x)$;

  \item находить пересечение графиков только «на глаз»: график полезен для оценки, точный ответ требует вычислений;

  \item забывать случай $a=0$: тогда квадратная функция вырождается в линейную.
\end{itemize}

% =========================================================
% БЛОК-СХЕМА 2
% =========================================================

\clearpage
\thispagestyle{fancy}

\begin{center}
{\Large\textcolor{darkaccent}{\textbf{Алгоритм: знаки коэффициентов параболы}}}
\end{center}

\vspace{9mm}

\begin{center}
\begin{tikzpicture}[node distance=9mm]

\node[flowstart] (start)
{Дан график $y=ax^2+bx+c$};

\node[flowproc,below=of start] (a)
{По направлению ветвей определите знак $a$};

\node[flowproc,below=of a] (xv)
{Определите знак $x_V$: вершина слева, на оси или справа от $Oy$};

\node[flowproc,below=of xv] (b)
{Используйте $b=-2ax_V$ и определите знак $b$};

\node[flowproc,below=of b] (c)
{По пересечению с осью $Oy$ определите знак $c$};

\node[flowcond,below=13mm of c] (check)
{Все три знака согласуются с графиком?};

\node[flowend,below=14mm of check] (end)
{Запишите знаки $a$, $b$, $c$};

\node[
  flowproc,
  right=10mm of check,
  text width=40mm
] (redo)
{Повторно проверьте вершину и пересечение с $Oy$};

\draw[flowarrow] (start) -- (a);
\draw[flowarrow] (a) -- (xv);
\draw[flowarrow] (xv) -- (b);
\draw[flowarrow] (b) -- (c);
\draw[flowarrow] (c) -- (check);

\draw[flowarrow]
  (check)
  --
  node[flowlabel,right]{да}
  (end);

\draw[flowarrow]
  (check.east)
  --
  node[flowlabel,above]{нет}
  (redo.west);

\draw[flowarrow]
  (redo.north)
  |-
  ([xshift=6mm]xv.east)
  --
  (xv.east);

\end{tikzpicture}
\end{center}

\clearpage

% =========================================================
% 5. ПАРАМЕТРЫ
% =========================================================

\section*{5. Параметр как движение графика}
\sectionline

Параметр --- это число, значение которого может меняться.
Полезно видеть, \emph{как именно} при этом меняется график.

\begin{tabularx}{\linewidth}{@{}>{\bfseries}p{0.24\linewidth}X@{}}
\toprule
Семейство & Геометрический смысл \\
\midrule
$y=m$
&
Горизонтальная прямая перемещается вверх и вниз.
\\

$y=mx$
&
Прямая проходит через $(0;0)$ и поворачивается при изменении $m$.
\\

$y=(m+2)x-3$
&
Меняется угловой коэффициент $m+2$, а точка пересечения с $Oy$ остаётся $(0;-3)$.
\\
\bottomrule
\end{tabularx}

Если требуется найти значение параметра, при котором два графика пересекаются
в заданной точке $(x_0;y_0)$, эта точка должна удовлетворять
\textbf{обеим} формулам.

Поэтому значения $x_0$ и $y_0$ подставляют в каждую зависимость
и проверяют совместность условий.

Для горизонтальной прямой $y=m$ число пересечений с параболой удобно оценивать
по вершине.

Например, для
\[
y=(x-h)^2+q
\]
имеем:

\[
\begin{array}{ccl}
m<q &\Longrightarrow& \text{общих точек нет},\\[1mm]
m=q &\Longrightarrow& \text{одна общая точка},\\[1mm]
m>q &\Longrightarrow& \text{две общие точки}.
\end{array}
\]

Для параболы с ветвями вниз ситуация меняется зеркально.

\textbf{ОДЗ.}
Все рассмотренные линейные и квадратичные функции являются многочленами,
поэтому
\[
x\in\mathbb{R}.
\]

\textbf{Пример проверки.}
Для
\[
y=4x+5
\]
точка $(2;1)$ не подходит, потому что
\[
4\cdot2+5=13\ne1.
\]

Следовательно, никакой выбор параметра во второй формуле не сделает
$(2;1)$ общей точкой, если первая функция остаётся
\[
y=4x+5.
\]

\subsection*{Мини-памятка}

\begin{center}
\begin{tabularx}{0.92\linewidth}{@{}X X@{}}
\toprule
\textbf{Линейная функция $y=kx+b$}
&
\textbf{Парабола $y=ax^2+bx+c$}
\\
\midrule
$k$ --- наклон
&
$a$ --- ветви
\\

$b$ --- пересечение с $Oy$
&
$b$ --- положение вершины через
$x_V=-\dfrac{b}{2a}$
\\

$f(x)=g(x)$ --- пересечение графиков
&
$c$ --- пересечение с $Oy$
\\
\bottomrule
\end{tabularx}
\end{center}

% =========================================================
% ТРЕНИРОВОЧНЫЙ БЛОК
% =========================================================

\clearpage

\section*{Тренировочный блок}
\sectionline

Выполняйте задания по порядку:
сложность постепенно возрастает.
Решения в этом блоке не приводятся.

\begin{enumerate}

  \item
  Для прямой
  \[
  y=-3x+4
  \]
  укажите знак углового коэффициента и координату точки пересечения с осью $Oy$.

  \item
  Какая из точек лежит на графике
  \[
  y=2x-5:
  \]
  \[
  A(1;-3),
  \qquad
  B(2;0),
  \qquad
  C(3;2)?
  \]

  \item
  Найдите точку пересечения прямых
  \[
  y=3x-4,
  \qquad
  y=-x+8.
  \]

  \item
  Для параболы
  \[
  y=ax^2+bx+c
  \]
  известно:

  ветви направлены вниз;

  вершина расположена справа от оси $Oy$;

  график пересекает ось $Oy$ выше оси $Ox$.

  Определите знаки $a$, $b$, $c$.

  \item
  Для параболы
  \[
  y=2x^2-8x+3
  \]
  найдите абсциссу вершины и укажите,
  справа или слева от оси $Oy$ она находится.

  \item
  Найдите точки пересечения прямой и параболы:
  \[
  y=x+1,
  \qquad
  y=x^2-3x+1.
  \]

  \item
  При каком значении параметра $m$ точка $(2;7)$ лежит на прямой
  \[
  y=mx+1?
  \]

  \item
  Для каких значений $m$ горизонтальная прямая
  \[
  y=m
  \]
  имеет с параболой
  \[
  y=(x-2)^2-1
  \]
  ровно одну общую точку?

  Укажите это значение $m$.

  \item
  Найдите значение параметра $m$, при котором прямые
  \[
  y=(m+2)x-3,
  \qquad
  y=x+2m
  \]
  пересекаются в точке с абсциссой
  \[
  x=1.
  \]

  Затем найдите ординату точки пересечения.

  \item
  Парабола
  \[
  y=ax^2+bx+c
  \]
  проходит через точку $(0;-2)$,
  её вершина имеет абсциссу
  \[
  x_V=3,
  \]
  а
  \[
  a=2.
  \]

  Найдите $b$ и $c$, затем запишите формулу параболы.

\end{enumerate}

% =========================================================
% ОТВЕТЫ
% =========================================================

\clearpage

\section*{Ответы к тренировочному блоку}
\sectionline

\begin{table}[h!]
\centering
\caption{Краткие ответы}

\begin{tabularx}{\linewidth}{@{}cX@{}}
\toprule
\textbf{№} & \textbf{Ответ} \\
\midrule

1
&
$k<0$;
точка пересечения с $Oy$: $(0;4)$.
\\

2
&
$A(1;-3)$.
\\

3
&
$(3;5)$.
\\

4
&
$a<0$,
$b>0$,
$c>0$.
\\

5
&
$x_V=2$;
вершина справа от оси $Oy$.
\\

6
&
$(0;1)$ и $(4;5)$.
\\

7
&
$m=3$.
\\

8
&
$m=-1$.
\\

9
&
$m=-2$;
точка пересечения $(1;-3)$.
\\

10
&
$b=-12$,
$c=-2$;
\[
y=2x^2-12x-2.
\]
\\

\bottomrule
\end{tabularx}
\end{table}

\vfill

\begin{center}
\textcolor{darkaccent}{\textbf{Формула для самопроверки:}}
\quad
\textit{сначала графический смысл, затем точный расчёт.}
\end{center}

\end{document}